\section{Configuración, Generación y Seguridad de la Base de Datos}

\subsection{Generación de un histórico realista}

Para evaluar el asistente en escenarios financieros verosímiles, necesitamos un 
conjunto de datos suficientemente rico. Para ello, el proyecto incluye un módulo 
de generación sintética, \texttt{seed.py}, que crea
un histórico de operaciones correspondiente a un periodo de 24 meses.

El algoritmo de generación es completamente determinista, inicializado con
\texttt{random.seed(42)}. Esto facilita tanto el desarrollo como la
repetición de las pruebas, ya que los resultados se pueden comparar sobre un
conjunto de datos conocido.

Los datos generados simulan patrones de consumo realistas que exigen distintos 
niveles de complejidad analítica por parte del asistente:

\begin{itemize}
    \item \textbf{Gastos recurrentes y variables:} se combinan cargos fijos y regualres 
    \item (como el aquiler, la luz o el agua) con transacciones con alta variabilidad 
    (compras en supermercados, repostaje en gasolineras, restaurantes, etc.).

    \item \textbf{Pagos con periodicidad anual:} se incluyen operaciones
    separadas por periodos de tiempo más amplios, como la renovación anual del
    seguro del coche. Esto permite comprobar el funcionamiento de las consultas
    que requieren aplicar filtros sobre periodos de tiempo extensos.

    \item \textbf{Cambios en los gastos:} e simulan subidas de cuota a mitad del histórico, 
    como incrementos documentados en el precio de plataformas de 
    \textit{streaming} para comprobar si el sistema es capaz de reconocer dicho incremento.
\end{itemize}

De esta forma, la base de datos proporciona un conjunto de pruebas más variado
que permite comprobar tanto consultas sencillas como aquellas que requieren
analizar información de diferentes periodos.

\subsection{Defensa en profundidad y aislamiento del LLM}

Una parte importante del sistema es controlar el acceso del modelo de lenguaje
a la base de datos. El modelo puede generar consultas para obtener información
a partir de las peticiones realizadas por el usuario, por lo que es necesario
evitar que dichas consultas puedan modificar los datos almacenados.

Para ello, el acceso a la base de datos, implementado en
\texttt{database.py}, incorpora varias medidas de protección independientes:

\begin{itemize}

    \item \textbf{Acceso de solo lectura:} cuando el modelo necesita consultar
    información del historial, el backend utiliza una conexión SQLite en modo
    solo lectura mediante \texttt{mode=ro}. De esta forma, las operaciones que
    intenten modificar la base de datos, como \texttt{DROP} o \texttt{UPDATE},
    son rechazadas por el propio motor de SQLite.

    \item \textbf{Validación de las consultas:} antes de ejecutar una consulta
    generada por el modelo, el sistema comprueba su contenido mediante una serie
    de expresiones regulares. Se bloquean palabras clave asociadas a operaciones
    de modificación de datos o de la estructura de la base de datos, como
    \texttt{INSERT}, \texttt{DELETE} o \texttt{ALTER}. Además, se comprueba que
    la consulta comience por \texttt{SELECT} o \texttt{WITH}.

    \item \textbf{Limitación del número de resultados:} las consultas que
    devuelven un número elevado de registros pueden generar una cantidad de
    información innecesaria para el modelo. Para evitarlo, el sistema establece
    un límite de 50 filas mediante la constante \texttt{MAX\_FILAS = 50}.
    Cuando una consulta devuelve más resultados, estos se truncan antes de
    incorporarse al contexto de la conversación.
    
\end{itemize}

Estas medidas permiten separar la generación de las consultas de su ejecución.
El modelo de lenguaje se encarga de construir la consulta a partir de la
petición del usuario, mientras que el backend controla qué consultas pueden
ejecutarse y con qué permisos se accede a la base de datos.